\chapter{The common source for matched spherical diagrams}
\label{app:morse-preliminaries}

A simple matched white sphere determines the complete augmented
geometry and the exact reduced finite source used in the following
volume theorems. This chapter records those common inputs once.
No terminal decoration or source-balance inequality is needed here.
The two faces adjacent to the marked strand edge are protected; after
sending their common contact to infinity they have curvature zero.
All other face variables are shared throughout the actual diagram.

\section{A simple sphere supplies the actual metric}\label{bl:a-simple-sphere-supplies-the-actual-metric}

\subsection{A combinatorial starting hypothesis}\label{bl:a-combinatorial-starting-hypothesis}

It is enough to start with an actual connected marked augmented diagram
whose white cellulation is a simple triangulated sphere with at least
four vertices. The matching pairs, cyclic orders and marked strand edge
are those of its projection/crossing-disk cut. A positive geometric
certificate need not be supplied independently.

Three steps turn the combinatorial sphere into a polyhedron: an
angle assignment, a planar development of its center triangles, and
the supporting-halfspace test. We use the following calibrated
triangle entropy. Remove one triangular
face \(abc\), prescribe positive radii at its three vertices, and let
\(v\) be the number of remaining interior vertices. The disk has
\(2v+1\) triangles. Give every interior vertex angle demand \(2\pi\),
and pool the three boundary vertices into one destination with demand
\(\pi\), retaining the separate boundary corners.

A nonempty interior subset \(S\) touches at least \(2|S|+1\) disk triangles.
To see this, puncture the sphere at a vertex of \(S\). The complementary
simple planar graph on the other interior vertices \(T\) and \(a,b,c\)
has at most \(2|T|+1\) bounded triangular faces. One is the removed
face, so at most \(2|T|\) of the disk's triangles avoid \(S\).
A proper subset containing the pooled boundary destination has strict
slack too: its complement consists only of interior vertices \(T\),
and at most \(2|T|-5\) triangles are wholly supported there when
\(|T|\ge3\), with none when \(|T|\le2\). The full set has equal
supply and demand. The finite transport Hall inequalities are therefore
strict for every nonempty proper destination subset. Subtracting a sufficiently
small positive allocation at every corner preserves feasibility of the
remaining transport system. Thus a strictly positive angle assignment
exists.

For an angle triple of sum \(\pi\), use

\begin{equation}\label{bl:G0}
\mathcal V(\theta)=2\sum_{j=1}^3
 [\Lambda(\theta_j/2)+\Lambda((\pi-\theta_j)/2)],
 \qquad \Lambda(t)=-\int_0^t\log|2\sin u|\,du .
\end{equation}

Its corner derivative is \(-\log\tan(\theta_j/2)\), and its one-corner endpoint value is zero. Maximize the sum of these calibrated triangle entropies, with the fixed
boundary linear terms \(\theta_f\log(1/r_f)\). Their derivative
\(-\log\tan(\theta/2)\) diverges positively at a new zero corner.
Moving toward the strictly positive assignment excludes a boundary
maximum. Strict concavity gives one interior maximum. Its multipliers
satisfy

\begin{equation}\label{bl:G1}
\tan(\theta_{tf}/2)=k_t l_f,\qquad
k_t=\left(\sum_{f<g\in t}l_fl_g\right)^{-1/2},
\qquad l_f=1/r_f>0.
\end{equation}

The boundary multipliers are the prescribed inverse radii, up to the
one common scale absorbed in the row multipliers. The three boundary
angle sums are positive and total \(\pi\), hence each is less than
\(\pi\).

For each triangle glue the Euclidean center triangle with side lengths
\(r_f+r_g\). Its angles are exactly \eqref{bl:G1}. The total angle at every
interior vertex is \(2\pi\); continuation around that vertex fixes
its center and rotates by \(2\pi\), so has zero translational as
well as rotational holonomy. These loops generate the punctured
disk's holonomy. The development is therefore globally defined. Its
boundary is the convex triangle with the prescribed three side lengths,
traversed once. For a point away from the developed edges, the number
of preimages is its boundary winding number: one inside and zero
outside. Local injectivity extends this across edges and vertices.
The development is an embedding.

Place the radius-\(r_f\) disk at each developed center. In one center
triangle the opposite center edge is covered by its two vertex disks,
both disjoint from the third disk. Their mutual contact lies strictly
outside the third disk, and its two contacts with those disks occur
on the other two center edges. Thus the third disk's sector stops
strictly before the opposite edge. The sectors assemble to full disks
at interior vertices and to the prescribed boundary-disk sectors.
This proves disjoint interiors and excludes unintended contacts,
including across separating triangles.

The circle through a center triangle's three tangency points is its
incircle: its contact distance along a center edge is exactly the
corresponding radius. It is orthogonal to the three vertex circles.
Its disk contains no other packing contact. On the omitted outer
triangle choose the opposite side of its contact circle; all interior
contacts lie in its bounded curvilinear gap. The resulting white and
shaded supporting halfspaces contain every contact, strictly unless
incident. Their exposed faces form the full closed labeled sphere,
so the convex hull has exactly those facets and no additional face.
Every ideal vertex has alternating white/shaded faces and successive
right angles. This constructs the required convex ideal right-angled
polyhedron. The proof did not use a height ordering or a flat plat.

Sending any chosen strand contact to infinity turns its adjacent white
circles into two parallel marked lines. Every other disk is finite in
their common strip. Thus the required marked complete packing follows
from the finite combinatorial starting hypothesis.

\section{The original link and a knot parity}\label{bl:the-original-link-and-a-knot-parity}

\subsection{From a matched sphere to an actual knot diagram}\label{bl:from-a-matched-sphere-to-an-actual-knot-diagram}

Conversely begin with any finite oriented simple triangulated sphere
with at least four vertices and a chosen perfect matching of its cubic
dual. Two matched triangles share one edge, have distinct third vertices,
and form an embedded quadrilateral with four distinct boundary vertices.
Choose disjoint small neighborhoods of the matched dual edges and their
endpoints. Contracting these disks gives a connected planar four-valent
graph. At each disk use the original ordered untwisted pairing
\((0,1),(2,3)\); join the outside terminals by the embedded nonmatching
dual edges. Put the standard perpendicular crossing circle in a disjoint
three-ball neighborhood of each disk. This constructs an augmented
link in \(S^3\), with exactly the prescribed cut graph and disk frames.

Choose a regular height after straightening each local box in its
orientation-preserving frame. The outside arcs give finitely many
ordinary cups and caps and no new crossing. This is the literal
general Morse diagram, not an assumed flat-plat presentation. A strand
edge may be chosen for the original reduced mark.

The graph whose vertices are the untwisted main circles and whose
edges are the crossing boxes is connected. Otherwise the union of
boxes and outside arcs would split the connected planar graph.
Choose a spanning tree of this graph and give exactly those boxes
odd parity. Switching one tree edge to the odd pairing joins two
different current circles. Induction along any order of the tree
edges gives one main component, with no internal extra component.
The odd signs are arbitrary. Subsequent full twists preserve this
pairing, and the even terminal fans preserve it as well. This proves
the combinatorial corollary without asserting that every arbitrary
tensor graph is a knot diagram.

\section{The complete augmented exterior}
\label{morse:complete-exterior}

Cut along the projection sphere. In each projection ball a crossing
circle semicircle is boundary parallel by its original half disk.
Drilling that semicircle and opening that same half disk cancels the
introduced handle. These operations occur in disjoint collars, so
opening all of them leaves one ball in each projection copy. The
boundary cellulation is the prescribed medial sphere. Thus cycles in the
matched graph create no unidentified handle or extra gluing.

Let \(Q\) be its positive ideal right-angled cut polyhedron. The even
shaded pairing fixes the crossing-circle vertex and sends
\(C_1\mapsto C_0,\ C_2\mapsto C_3\). An odd pairing exchanges the
projection copies and the two punctures, using the literal signed
half-twist collar. The two odd signs have the same corner permutation
but different primitive paths. White faces pair by their original
projection labels.

At every ideal vertex, the link of \(Q\) is a rectangle. Normalize its
opposite shaded sides to length one. Each shaded triangle map then
preserves the canonical horocycles at all three corners, including the
odd corner permutation. Each alternating white/shaded edge cycle has
four right-angle pieces. It fixes both ideal endpoints, preserves the
chosen horospheres there, and rotates by \(2\pi\), so its whole holonomy
is the identity. The ideal links are the actual link's torus links and
have no cone or dilation defect. Their horoball ends are complete. Hence
\begin{equation}\label{mcg:eq-4}
\operatorname{Vol}(Y_D)=2\operatorname{Vol}(Q),
       \qquad \Phi_D=\frac{\operatorname{Vol}(Y_D)}{2\pi}.
\end{equation} Here \(Y_D\) is the same augmented exterior, with the recorded half
twists and peripheral bases.

The construction uses the original sphere projection and disjoint
crossing-disk collars. In particular it applies before any terminal fan
is inserted and does not require a flat-plat presentation.

\section{The calibrated triangle action}\label{bl:an-exact-integral-for-the-source-normalized-triangle-action}

Let \(F(x,b,y)\) be the complete triangle action with the original
determinant, beta and half-forced-product normalization. Put
\(A=x+y\) and \(P=xy\). Its exact
entropy definition is

\begin{equation}\label{bl:F8}
F(x,b,y)=\max_{\alpha_x+\alpha_b+\alpha_y=1/2}
\left\{\sum_j f(\alpha_j)+\alpha_x\log x+
\alpha_b\log b+\alpha_y\log y\right\},
\end{equation}

where \(f'(t)=\log\cot(\pi t)\), \(f(0)=f(1/2)=0\), and zero
weights force zero allocation. The complete stationary middle load is

\begin{equation}\label{bl:F9}
\frac{\partial F}{\partial\log b}
=\frac1\pi\arctan\frac{b}{\sqrt{Ab+P}}.
\end{equation}

The entropy derivatives give these identities with their displayed
absolute constant; the same normalization is used in
\cite[Section~5]{CasimirSources} and \cite[Section~1]{MarkedFillings}.

\section{The physical face adapter}
\label{morse:face-adapter}

Fix \(N\ge3\), \(d=N-1\), \(q=e^{\pi i/N}\), \(Q=q^2\). Let \(H_f(x,b,c,y)\) be the original homogeneous degree-\(d\) face polynomial of an arbitrary full coupon \(f(\Omega)\). This includes actual signed crossing powers and cleared resolvent numerators, with all repeated-node data.

The finite Fourier face adapter of \cite[equation~(4.5)]{Finite} is

\begin{equation}\label{bl:S5}
f(\Omega)_{r,s;u,v}
 =N^{-1}q^{-d(r-u)}c^{-d}H_f(x,b,c,y)
 =N^{-1}\frac{H_f(x,b,c,y)}{(bc)^{d/2}} .
\end{equation}

The last equality uses the coherent integer face-log lifts of that physical state. No head is independently reduced modulo \(N\). The physical Casimir lowering arrows vanish at their genuine endpoints; hence a modular output with a wrong unwrapped total has zero coefficient before it is interpreted as a physical state. This is valid for every polynomial in the full Casimir, not only for the complete circle.

Let \(M_f\) count middle incidences and \(T_f\) outside minus
inside Morse turns in a face, temporarily restoring the cap at the
marked edge. Rounding the boxes in their actual frames, each middle
passage contributes \(+\pi\) to boundary turning, while the outer pair
contributes zero. Inside and outside turns contribute \(+\pi\) and
\(-\pi\). The winding is one in a bounded face and minus one in the
exterior face \(a\). Consequently
\begin{equation}\label{bl:S8}
 M_f-T_f=2\quad(f\ne a),\qquad M_a-T_a=-2.
\end{equation}
The equality uses the actual region boundary, so adding a cancelling
Morse cup--cap pair leaves it unchanged.

\section{The physical mark and the full finite source}\label{bl:the-physical-mark-measure-and-normals}

Choose a strand edge as the mark, with adjacent protected faces
\(a,b_*\). The pivotal self-duality transports its reduced cut through
the actual Morse presentation; it does not replace an internal marked
row by an ordinary cap. The regular generic pivot \((-1)^HK\), where
\(H\) is the integer-weight operator, specializes to \(K^{1-N}\).
The transported cut and literal duality frames are those of
\cite[Section~9]{Finite}.

In literal left-to-right physical order the cup, cap and mark are

\begin{equation}\label{bl:S12}
p_{ij}=Nq^{-d^2+di}\mathbf1_{i+j=d},\qquad
 \ell_{ij}=N^{-1}q^{di}\mathbf1_{i+j=d},\qquad
 m_{ij}=N^{-2}q^{dj}.
\end{equation}

The marked row is initially unrestricted. The whole-network degree conservation, after the other actual caps act, forces its two heads to sum to \(d\). Only then is its ratio to the ordinary cap

\begin{equation}\label{bl:S13}
\frac m\ell=N^{-1}(L_{b_*}/L_a)^d .
\end{equation}

The cup and cap factors have cancelling powers of \(N\) and cancelling
fixed \(q\)-constants. For a diagram with \(R\) boxes,
\eqref{bl:S5}, \eqref{bl:S8} and \eqref{bl:S13} make its full marked summand

\begin{equation}\label{bl:S14}
N^{-R-1}\,
 \frac{\prod_{e=1}^{R}H_{f_e}(L_{x_e},L_{b_e},L_{c_e},L_{y_e})}
      {\prod_{f\ne a,b_*}L_f^d}.
\end{equation}

The integer conservation law is a face-potential law on the sphere.
Integrate head differences along dual paths: zero divergence makes
the result path independent, up to one common constant. This gives a
primitive lattice, unchanged by bending an arc. There are \(R+2\)
faces and hence \(R+1\) independent \(N\)-point grids after fixing
\(L_a\). Their count cancels \(N^{-R-1}\) in \eqref{bl:S14}.
There is no residual dimension factor in this conversion.

The numerator is homogeneous of degree \(dR\), with nonnegative
exponents and arbitrary complex coefficients. The following finite
orthogonality statement gives exactly the required coefficient.

\begin{proposition}[The unique monomial selected by the face grids]
\label{mcg:face-grid-selection}
Let \(N\ge2\), \(d=N-1\), and let \(P\) be a polynomial in
\(R+2\) face variables, of total degree at most \(Rd\). Distinguish
two variables \(L_u,L_v\) and call the remaining \(R\) variables
ordinary. Put \(\omega=e^{2\pi i/N}\). For arbitrary nonzero
\(\xi_f\in\mathbb C\), \(f\ne u\), set
\(L_u=1\) and \(L_f=\xi_f\omega^{p_f}\). Then
\[
 \frac1{N^{R+1}}
 \sum_{\boldsymbol p\in\{0,\ldots,N-1\}^{R+1}}
 \frac{P(L)}{\prod_{f\ne u,v}L_f^d}
 =
 [\prod_{f\ne u,v}L_f^d]P\big|_{L_u=L_v=0}.
\]
In particular, translating any of the free grids by a nonzero factor
does not change this average.
\end{proposition}

\begin{proof}
By linearity it suffices to take a monomial with nonnegative exponents
\(m_f\). The finite geometric sum in each ordinary variable vanishes
unless \(m_f\equiv d\pmod N\); the grid at \(v\) also requires
\(m_v\equiv0\pmod N\). Every ordinary exponent that survives is
at least \(d\). Their sum is therefore at least \(Rd\), which
exhausts the assumed total-degree bound. Hence every ordinary exponent
equals \(d\) and \(m_u=m_v=0\). The remaining multiplier is
\(\prod_{f\ne u,v}\xi_f^{m_f-d}\xi_v^{m_v}=1\).
This monomial contributes its coefficient; every other monomial
contributes zero. No sign or positivity condition on the coefficients
was used.
\end{proof}

Each original \(H_e\) is homogeneous of degree \(d\), so their
product has degree exactly \(Rd\). At even levels the physical grids
are translated by \(\xi_f=q^{t_f}\). The forced exponents in the
proposition show that their extra character is one. Thus the original
marked scalar is

\begin{equation}\label{bl:S15}
I_N=\left[\prod_{f\ne a,b_*}L_f^d\right]
 \left.\prod_{e=1}^{R}H_{f_e}
       (L_{x_e},L_{b_e},L_{c_e},L_{y_e})
 \right|_{L_a=L_{b_*}=0}.
\end{equation}

This holds for arbitrary local full polynomials and their spectral parameters. In particular use

\begin{equation}\label{bl:S16}
V_N(\sigma;z)=
  \frac{\chi_N(\sigma)-\chi_N(z)}{\sigma-z},\qquad
 \chi_N(z)=2T_N(z/2)+2,
\end{equation}

as a polynomial before setting \(z=\Omega\), where
\(T_N(\cos t)=\cos Nt\). It is the actual cleared resolvent at the root
and retains the first derivative at every double Casimir node. No
value-only interpolation or paired-channel retraction is used.

If an ordinary generic closure is used instead, the full open regular network is a scalar endomorphism of the generic highest-weight object \(V_d\). Its ordinary closure is \((-1)^d[N]_q\) times that scalar. The one genuine normal at the root recovers \eqref{bl:S15}. No additional normal occurs at a crossing circle or at a separator. This uses the regular physical lift, not a derivative of an arbitrary fixed-root identity.

\subsection{The original signed Fourier/Hermite period}\label{bl:the-original-signed-fourierhermite-period}

Put the literal nonzero fixed signed braid power \(R^{n_e}\) at each box, with \(n_e=\epsilon_e-2p_e\) in its own original disk framing. Define the whole spectral numerator

\begin{equation}\label{bl:S17}
S_N(\boldsymbol\sigma)=
 \left[\prod_{f\ne a,b_*}L_f^d\right]
 \left.\prod_{e=1}^{R}H_{V_N(\sigma_e)}\right|_{L_a=L_{b_*}=0}.
\end{equation}

Apply the source-independent local full-Hermite crossing functional
separately to each box, as in \cite[Corollary~9.2]{Finite}.
Multilinearity of the original ordered contraction gives

\begin{equation}\label{bl:S18}
\begin{split}
 I_{N,\boldsymbol n}={}&
 q^{\frac{d^2}{4}\sum_e n_e}(-i)^{R}
 \sum_{\boldsymbol\ell\in(\mathbb Z\setminus\{0\})^{R}}
 \left(\prod_e\ell_e\right)(-1)^{(N+1)\sum_e\ell_e}\\
 &\times\int_{[-1/2,1/2]^{R}}
 \left(\prod_e\sin(2\pi\delta_e)\right)
 S_N(-2\cos(2\pi\boldsymbol\delta))\\
 &\qquad\times
 \exp\!\left\{N\sum_e[\pi i n_e\delta_e^2+2\pi i\ell_e\delta_e]\right\}
 d\boldsymbol\delta .
 \end{split}
\end{equation}

This is the original whole Hermite functional applied to the general
Morse source. The Cartan scalar, charge weights, parity and endpoint
sine factors have not changed. The original matrix identity excludes
charge zero. Three integrations by parts on whole intervals give
absolute summability, retaining the second-derivative endpoint currents.
No constant-headcount assumption occurs.

The exterior framing/component character remains separately recorded. At complete main color it has modulus one. The full one-cut modified convention also retains its specified polynomial-level factor. These source conversions do not change the exponential rate. The invariant in \eqref{bl:S18} belongs to the specified filled diagram, with its original component and framing convention.

